% Source-derived chapter 05: Parabolic support theorem
% Original source: pw-conjecture-gl-n
\chapter{Parabolic support theorem}
\label{pw-conjecture-gl-n-chapter-05}
\source{pw-conjecture-gl-n}

\begin{remark}[Source section]
\label{pw-conjecture-gl-n-chapter-05-source}
\source{pw-conjecture-gl-n}
This chapter records the source section; no Lean declaration has yet been checked for this material.
\notready
\end{remark}

% The source section title was: Parabolic support theorem
In the last section, we prove Theorem \ref{thm3.2}. This is the parabolic version of the Chaudouard--Laumon support theorem \cite{CL}. For the proof, we ultimately reduce it to a relative dimension bound which we establish in Section \ref{sec4.4}.

\subsection{Review of support theorems}\label{Sec4.1}
\source{pw-conjecture-gl-n}

We start with a review of support theorems for Hitchin systems.

For a proper morphism $f: X\to Y$ with $X,Y$ nonsingular Deligne--Mumford stacks, the decomposition theorem \cite{BBD} yields
\[
Rf_* \BQ_X \simeq \bigoplus_i {^\mathfrak{p}\CH^i(Rf_* \BQ_X)[-i]}
\]
with ${^\mathfrak{p}\CH^i}(Rf_* \BQ_X)$ semisimple perverse sheaves on $Y$; we say that a closed $Z \subset Y$ is a support of $f$ if it is a support of a simple summand of some ${^\mathfrak{p}\CH^i}(Rf_* \BQ_X)$. A particularly interesting case is that $f$ has full support, that is, $Y$ is the only support of $f$; in this case the cohomology of any closed fiber of $f$ is governed by the nonsingular fibers.

The study of supports for Hitchin systems was initiated by B.C. Ng\^o and is crucial in his proof of the fundamental lemma of the Langlands program \cite{Ngo}. He determines all the supports for the Hitchin system (including the $\CL$-twisted cases) after restricting to the elliptic locus of the Hitchin base; that is the subset formed by integral spectral curves.

After Ng\^o's work, Chaudouard--Laumon \cite{CL} observed that, if we consider the moduli space of $\CL$-twisted $\mathrm{GL}_n$ stable Higgs bundles with $\mathrm{deg}(\CL) >0$, then Ng\^o's support theorem can be extended to the total Hitchin base; in particular they showed that the $\CL$-twisted $\mathrm{GL}_n$ Hitchin system has full support. Chaudouard--Laumon's idea was extended to the $\mathrm{SL}_n$ case \cite{dC_SL}, the endoscopic moduli spaces \cite{MS}, and singular cases involving strictly semistable Higgs bundles \cite{MS2,MS3}. See also \cite{dCHM2,MM} concerning the supports over the open subset of reduced spectral curves for the untwisted (\emph{i.e.} $\CL = 0$) Hitchin system.

The idea of Chaudouard--Laumon \cite{CL} is to show that each support of the $\CL$-twisted Hitchin system has generic point lying in the elliptic locus; this is achieved by combining two constraints: (I) the support inequality for a weak abelian fibration, and (II) $\delta$-regularity for spectral curves. We describe them in more detail.
\begin{enumerate}
    \item[(I)] \emph{(Support inequaltiy)} The Hitchin system admits the structure of a weak abelian fibration, that is, there is a commutative group scheme $P$ over the Hitchin base acting on the moduli space which satisfies certain properties. Then an argument generalizing the Goresky--MacPherson inequality leads to a codimension estimate for the supports. More concretely, it says that any support $Z$ has codimension bounded above by the $\delta$-function of the group $P$. This part was already carried out in Ng\^o \cite[Section 7]{Ngo}; see \cite[Section 1]{MS2} for a summary.
    \item[(II)] \emph{($\delta$-regularity)} In the case of type $A$ and for the stable locus, the group scheme is obtained from the multi-degree 0 relative Picard (or, for $\mathrm{SL}_n$, Prym) variety associated with the family of spectral curves. Then we have the Severi inequality, referred to as \emph{$\delta$-regularity}, for the spectral curves. We refer to \cite[Theorem 5.4.4]{dC_SL} for the precise statement.
\end{enumerate}

Using (I) and (II), one can deduce that no support is allowed to have generic point lying outside the open subset of integral curves; this was explained in \cite[Section 11]{CL} for $\mathrm{GL}_n$, in \cite[Section 6.2]{dC_SL} for $\mathrm{SL}_n$, and in \cite[Section 4.5]{MS2} for a general complete linear system in a del Pezzo surface. The argument is to combine the two inequalities above to deduce a numerical contradiction if there is a support that appears outside the elliptic locus.

%We note that, for the Chaudouard--Laumon argument to work, we have to work with semistable $\CL$-twisted Higgs bundles with $\mathrm{deg}(\CL) > 0$.

\subsection{Parabolic Hitchin systems}
Now we focus on the $\CL$-twisted parabolic Hitchin system (\ref{parah}). Let $\widehat{A}^{\mathrm{ell}} \subset \widehat{A}$ be the open subset parameterizing integral spectral curves (the elliptic locus). Following Ng\^o's method, Yun \cite{Yun2,Yun3} proved a parabolic support theorem over $C \times \widehat{A}^{\mathrm{ell}}$ and determined all the supports of (\ref{parah}) in $C \times \widehat{A}^\mathrm{ell}$. More precisely, by \cite[Section 2]{Yun3} any strict subset of $C \times \widehat{A}^{\mathrm{ell}}$, which is a support of $\widehat{h}^{\mathrm{par}}|_{C \times \widehat{A}^{\mathrm{ell}}}$, is a component of the \emph{endoscopic loci} governed by the endoscopic theory of $G$. As we only consider the special case $G = \mathrm{PGL}_n$, there are no nontrivial endoscopic loci and the restricted Hitchin map $\widehat{h}^{\mathrm{par}}|_{C \times \widehat{A}^{\mathrm{ell}}}$ has full support.

To prove Theorem \ref{thm3.2}, it suffices to show that there is no support of (\ref{parah}) whose generic point lying outside $C \times \widehat{A}^{\mathrm{ell}}$.

Since stability for a $\mathrm{PGL}_n$ Higgs bundle is described by its corresponding vector bundle (see (\ref{233})), it is more convenient to work with the $\mathrm{SL}_n$ moduli spaces. We consider the following Cartesian diagram
\begin{equation}\label{diag321}
\source{pw-conjecture-gl-n}
\text{[Diagram omitted; see source PDF.]}
\end{equation}
Here the horizontal maps are given by the natural quotient maps by $\Gamma = \mathrm{Pic}^0(C)[n]$. To describe the map $\pi'$ at the left-side of the diagram (see \cite[Example 2.2.5]{Yun1}),  we recall that $\widecheck{M}$ parameterizes traceless Higgs bundles with fixed determinant
\[
(\CE, \theta) , \quad \theta: \CE \to \CE\otimes \Omega_\CL, \quad \mathrm{rk}(\CE) = n, ~~\mathrm{det}(\CE) \simeq \CN,~~ \mathrm{trace}(\theta) = 0
\]
with respect to slope stability. Similarly $\widecheck{M}^{\mathrm{par}}$ parameterizes
\[
(x, \CE = \CE_0 \supset \CE_1 \supset \cdots \supset \CE_n =  \CE_0(-x), \theta)
\]
where $x\in C$, $(\CE, \theta) \in \widecheck{M}$, each $\CE_i$ in the flag is of rank $n$ with $\CE_i/\CE_{i+1}$ a length 1 skyscraper sheaf supported at $x$, the Higgs field preserves the flag $\theta(\CE_i) \subset \CE_i\otimes \Omega_\CL$, and the map $\pi'$ in the diagram is the forgetful map. We note that the stability condition for a parabolic Higgs bundle is determined by the stability of the underlying Higgs bundle.

Following \cite[Example 2.2.5]{Yun1} we may also describe $\widecheck{M}^{\mathrm{par}}$ via spectral curves. If we present a point in $\widecheck{M}$ as a 1-dimensional sheaf $\CF_\alpha$ supported on a spectral curve $C_\alpha \subset \mathrm{Tot}(\Omega_\CL)$ with $p_\alpha: C_\alpha \to C$ the projection, then for any $x\in C$, a closed point in the fiber $\pi'^{-1}(x, \CF_\alpha)$ is represented by
\[
\CF_\alpha = \CF_0 \supset \CF_1 \supset \cdots \supset \CF_n = \CF_0\otimes p_\alpha^* \CO_C(-x), \quad \mathrm{lengh}(\CF_i/\CF_{i+1}) = 1;
\]
see \cite[(2.4)]{Yun1}.

Now we consider the $\mathrm{SL}_n$ parabolic Hitchin system
\begin{equation}\label{SLh}
\source{pw-conjecture-gl-n}
\widecheck{h}^{\mathrm{par}}: \widecheck{M}^{\mathrm{par}} \to C\times \widecheck{M} \to C\times \widehat{A}.
\end{equation}
From the diagram (\ref{diag321}), it recovers the $\mathrm{PGL}_n$ Hitchin system $\widehat{h}^{\mathrm{par}}: \widehat{M}^{\mathrm{par}} \to C \times \widehat{A}$ by taking the quotient of the source by $\Gamma$. In particular, we have
\[
R{\widehat{h}^{\mathrm{par}}}_* ~\BQ_{\widehat{M}^{\mathrm{par}}} = \left(R{\widecheck{h}^{\mathrm{par}}}_* ~\BQ_{\widecheck{M}^{\mathrm{par}}}\right)^\Gamma \in D_c^b(C\times \widehat{A}).
\]
Hence in order to prove Theorem \ref{thm3.2}, it suffices to show that there is no support of (\ref{SLh}) with generic point lying outside $C \times \widehat{A}^{\mathrm{ell}}$. In the next two sections, we adapt the strategy of Chaudouard--Laumon \cite{CL} and de Cataldo \cite{dC_SL} (for the $\mathrm{SL}_n$-version) to this parabolic setting, and verify the parabolic analogs of (I) and (II) of Section \ref{Sec4.1}.




\subsection{Weak abelian fibrations and $\delta$-regularity}

A general approach for proving support theorems was given in \cite[Theorem 1.8]{MS2}. We apply it here for the parabolic Hitchin system (\ref{SLh}). In this section, we explain how existing results allow us to reduce Theorem \ref{thm3.2} to a relative dimension bound (\ref{relative}). Then in the next section we prove this bound.

For the reader's convenience, we first recall some of the necessary ingredients. As mentioned in (I) of Section \ref{Sec4.1}, the notion of \emph{weak abelian fibration} introduced by Ng\^o \cite{Ngo} plays a key role. A weak abelian fibration is a triple $(\CP, \CM, \CA)$ with morphisms
\[
f: \CM\to \CA, \quad g: \CP \to \CA,
\]
such that all $\CP, \CM, \CA$ are nonsingular, $f$ is proper, $g: \CP \to \CA$ is a smooth group scheme, and $\CP$ acts on $\CM$ relatively over $\CA$; they satisfy the following conditions:
\begin{enumerate}
    \item[(i)] every closed fiber of the $\CP \to \CA$ has dimension $= \mathrm{dim}\CM - \mathrm{dim}\CA$,
    \item[(ii)] the action of $\CP$ on $\CM$ has affine stabilizers, and
    \item[(iii)] the Tate module associated with $\CP \to \CA$ is polarizable.
\end{enumerate}
We refer to \cite[Section 1.1]{MS2} for more details.

Now we show that (\ref{SLh}) is naturally enhanced into a weak abelian fibration. We first construct the $(C\times \widehat{A})$-group scheme $P$ which acts on $\widecheck{M}^{\mathrm{par}}$. Recall that in \cite{dC_SL} de Cataldo showed that $\widecheck{M}$ admits a weak abelian fibration structure $(\widecheck{P}, \widecheck{M}, \widehat{A})$. Here the $\widehat{A}$-group scheme $\widecheck{P}$ which acts on $\widecheck{M}$ is given by the identity component of the relative Prym variety associated with the spectral curves \cite[(43)]{dC_SL}. For a spectral curve $p_\alpha: C_\alpha \to C$ with a Higgs bundle given by $\CF_\alpha$ supported on $C_\alpha$, the Prym action is induced by tensor product
\[
\CQ \cdot \CF_{\alpha} = \CQ \otimes \CF_\alpha, \quad \CQ \in \mathrm{Prym}^0(C_\alpha/C) = \widecheck{P}_{[C_\alpha]},~ \quad [C_\alpha] \in \widehat{A}.
\]
We denote by $P$ the $(C\times \widehat{A})$-group scheme obtained by the pullback of $\widecheck{P}$. The $P$-action on $C\times \widecheck{M}$ can be lifted to $\widecheck{M}^{\mathrm{par}}$:
\[
\CQ \cdot (x, \CF_\alpha = \CF_0 \supset \CF_1 \supset \cdots \supset \CF_n) =  (x, \CQ\otimes\CF_\alpha = \CQ \otimes \CF_0 \supset \CQ\otimes\CF_1 \supset \cdots \supset \CQ \otimes \CF_n),
\]
since the stability condition does not rely on the flag.


\begin{prop}
\source{pw-conjecture-gl-n}
\notready\label{prop4.1}
The triple $(P, \widecheck{M}^{\mathrm{par}}, C\times \widehat{A})$ forms a weak abelian fibration, \emph{i.e.},  it satisfies (i,ii,iii) above or of \cite[Section 1.1]{MS2}.
\end{prop}

\begin{proof}
    (i) only concerns the dimension of $\widecheck{M}$ which is clear from (\ref{diag0}), and (iii) only depends on the group scheme $P$ which follows from the corresponding property for $\widecheck{P}$ proved in \cite[Theorem 4.7.2]{dC_SL}. To prove (ii): since the $P$-action on $\widecheck{M}^{\mathrm{par}}$ is a lifting of the $P$-action on $C\times \widecheck{M}$, and the latter has affine stabilizers by \cite{dC_SL}, we obtain that the $P$-stabilizer for any point $z \in \widecheck{M}^{\mathrm{par}}$ is contained in the (affine) $P$-stabilizer for $\pi(z) \in C \times \widecheck{M}$. This proves (ii).
\end{proof}


%We obtain immediately the $\delta$-regularity ((II) of Section \ref{Sec4.1}) of the weak abelian fibration $(P, \widecheck{M}, C\times \widehat{A})$. Indeed, this is a property only dependent on $P$, which follows from the corresponding property established for $\widecheck{P}$ \cite[corollary 5.4.4]{dC_SL}.

Next, we consider the $\delta$-function on the base $C\times \widehat{A}$. The group scheme $P$ endows $C\times \widehat{A}$ with an upper semi-continuous function
\[
\delta: C\times \widehat{A} \to \BN
\]
calculating the dimension of the affine part of the commutative group scheme given by each closed fiber. For a closed subset $Z\subset C\times \widehat{A}$, we define $\delta_Z$ to be the minimal value of the function $\delta$ on $Z$.

We say that the weak abelian fibration  $(P, \widecheck{M}^{\mathrm{par}}, C\times \widehat{A})$ given by Proposition \ref{prop4.1} satisfies the \emph{support inequality} ((I) of Section \ref{Sec4.1}), if the inequality
\begin{equation}\label{delta_ineq}
\source{pw-conjecture-gl-n}
\mathrm{codim}_{C\times \widehat{A}}Z \leq \delta_Z
\end{equation}
holds for any irreducible support $Z \subset C\times \widehat{A}$ associated with (\ref{SLh}). Furthermore, by \cite[Theorem 1.8]{MS2}, the support inequality (\ref{delta_ineq}) follows from the relative dimension bound
\begin{equation}\label{relative}
\source{pw-conjecture-gl-n}
   \tau_{> 2\mathbf{d}}\left( R\widecheck{h}^{\mathrm{par}}_* \BQ_{\widecheck{M}^{\mathrm{par}}}\right) = 0, \quad \mathbf{d}: = \mathrm{dim} \widecheck{M}^{\mathrm{par}} - \mathrm{dim} (C\times \widehat{A})
\end{equation}
with $\tau_{>*}$ the standard truncation functor.

Once we have (\ref{relative}),
we can combine the support inequality (\ref{delta_ineq}) with $\delta$-regularity of
$\widecheck{P}$ \cite[corollary 5.4.4]{dC_SL} to prove
Theorem \ref{thm3.2} as follows.  If $Z \subset C \times \widehat{A}$ is an irreducible support of $\widecheck{h}^{\mathrm{par}}$, the projection $W = \mathrm{pr}_{\widehat{A}}(Z) \subset \hat{A}$ also satisfies the support inequality
\[
\mathrm{codim}_{\widehat{A}} W \leq \mathrm{codim}_{C\times \widehat{A}}Z \leq \delta_Z = \delta_W.
\]
Here we used in the second inequality the support inequality (\ref{delta_ineq}) for $\widecheck{h}^{\mathrm{par}}$, and used in the last equality that the $\delta$-function on $C\times \widehat{A}$ is pulled back from $\widehat{A}$. The argument of \cite[Section 6.2]{dC_SL} then shows that the generic point of $W$ lies in $\widehat{A}^{\mathrm{ell}}$. Hence the generic point of $Z$ lies in $C\times \widehat{A}^{\mathrm{ell}}$ as we desired.


%This is because the argument there only depends on $(\widecheck{P}, \widehat{A})$, which automatically works for $(P, C\times \widehat{A})$.


In conclusion, this reduces Theorem \ref{thm3.2} to (\ref{relative}).



\subsection{Proof of the relative dimension bound (\ref{relative})}\label{sec4.4}
\source{pw-conjecture-gl-n}

We need to show that each closed fiber of the morphism
$$\widecheck{h}^{\mathrm{par}}: \widecheck{M}^{\mathrm{par}} \rightarrow C\times \widehat{A}$$
has dimension
\[
\mathbf{d} = \mathrm{dim}\widecheck{M}^{\mathrm{par}} - \mathrm{dim} (C\times \widehat{A}) =    \frac{n(n-1)}{2}\deg(\CL) + (n^2-1)(g-1).\]
Here the last equation is obtained by the dimension formula of \cite[Proposition 2.4.6]{dC_SL} directly:
\[
\mathbf{d} =  \mathrm{dim}\widecheck{M} - \mathrm{dim}\widehat{A} = \frac{n(n-1)}{2} \deg(\CL) + (n^2-1)(g-1).
\]

%$$\dim \widetilde{\CM_{L}} - \dim \CA_{L}- 1 = \frac{1}{2}(n-1)(n)\deg L + (n-1)(g-1).$$


First,  we note that the morphism $\widecheck{h}^{\mathrm{par}}$ is surjective since the usual Hitchin map $\widecheck{h}: \widecheck{M} \to \widehat{A}$ is surjective. Furthermore, as $\widecheck{h}^{\mathrm{par}}$ is equivariant with respect to the scaling $\BG_m$-action on the Higgs fields, it suffices to bound from above the dimension of the fiber over $(x, 0) \in C\times \widehat{A}$ for each point $x \in C$ by $\mathbf{d}$, since upper semicontinuity will force all other fibers to have the same dimension upper-bound.\footnote{Furthermore, we actually show that each fiber has the same dimension. But the dimension bound is enough for our purpose.}  We fix the point $p$ from now on.

Using the assumption $\deg(\CL)>2g$, we may express $\CL$ as $\CO_C(D)$ with $D = x_0 + x_1 + \dots + x_t$ an effective reduced divisor containing $x = x_0$. In particular, a Higgs bundle $(\CE, \theta: \CE \to \CE\otimes \Omega_\CL)$ can be viewed as a meromorphic (un-twisted) Higgs bundles with at most simple poles along $D$:
\[
(\CE, \theta), \quad \theta: \CE \to \CE \otimes \Omega_C^1(D).
\]
We will control the fiber dimension using an analogous bound for the nilpotent cone for \emph{strongly parabolic} Higgs bundles for the pair $(C,D)$.

Recall that, given $D$ as above, a strongly parabolic Higgs bundle consists of an $\mathrm{SL}_n$ Higgs bundle of degree $d$:
\[
(\CE, \theta), \quad \theta: \CE\to \CE \otimes \Omega^1_C(D),~~~\mathrm{rk}(\CE) =n,~~\mathrm{det}(\CE)\simeq \CN, ~~\mathrm{trace}(\theta) = 0,
\]
along with a flag at each point $x_i \in D$:
\[
\CE|_{x_i}=F_{x_i,0}  \supset F_{x_i,1} \supset F_{x_i,2}, \supset \cdots \supset F_{x_i,n} = 0, \quad \mathrm{dim} F_{x_i,j}/F_{x-i, j+1} = 1,
\]
such that the Higgs field satisfies $\theta(F_{x_i, j}) \subset F_{x_i, j+1}$.\footnote{One may compare this with the (non-strong) parabolic condition $\theta(F_{x,j}) \subset F_{x,j}$ we used earlier in this paper for the global Springer theory.}


%a vector bundle $\CE$ on $C$ with fixed determinant $\mathrm{det}(\CE) \simeq \CN$, along with reduction of structure group to the Borel subgroup $B$ at each point of $D$, equipped with a meromorphic Higgs field $\phi \in H^0(C, \mathfrak{g}\otimes K_C(D))$ such that the residue of $\phi$ at each point of $D$ lies in the $\mathfrak{n}$.  In terms of a flag of subspaces $F_n \subset F_{n-1} \subset \dots$, this last condition translates to requiring $\phi(F_j) \subset F_{j+1}$ for each $j$.


Let $\widecheck{M}^{\mathrm{spar}}(D)$ denote the moduli space of strongly parabolic Higgs bundles associated with the pair $(C,D)$, such that the underlying twisted Higgs bundle is stable. Let
$$A(D) = \bigoplus_{i=2}^{n} H^0(C, \Omega_\CL^{\otimes i}(-D)) \subset \widehat{A}$$
denote the linear subspace consisting of sections which vanish along the points of $D$. There is a Hitchin map for the strongly parabolic Higgs moduli space
$$\widecheck{h}^{\mathrm{spar}}:\widecheck{M}^{\mathrm{spar}}(D) \rightarrow A(D)$$
which is proper, surjective, and Lagrangian with respect to a natural holomorphic symplectic form \cite{Faltings}. In particular, the dimension of the zero fiber is given by
\begin{align*}
\dim \widecheck{M}^{\mathrm{spar}}(D)_0 = \dim A(D) &= \sum_{j=2}^{n} \left((j-1)(\deg(\Omega_L)) +g-1\right)\\
&= \frac{n(n-1)}{2}\deg (\CL) + (n^2-1)(g-1).
\end{align*}
See also \cite[Theorem 6.9]{SWW} for a direct proof of the above dimension formula.\footnote{The formula above has $g$ less than the dimension formula obtained in \cite[Theorem 6.9]{SWW}, since we consider the $\mathrm{SL}_n$ case where we fixed the determinant on $C$.}

We now consider the following diagram relating the two types of parabolic Higgs moduli spaces $\widecheck{M}^{\mathrm{spar}}(D)$ and
$\widecheck{M}^{\mathrm{par}}$:

\begin{equation}\label{diagram}
\source{pw-conjecture-gl-n}
\text{[Diagram omitted; see source PDF.]}
\end{equation}
where all squares are Cartesian.

We first claim that the natural map $q$ (sending a strong parabolic Higgs bundle to a parabolic Higgs bundle) is surjective. Indeed, given a point $z \in \widecheck{M}^{\mathrm{par}}$ lying over $C\times {A}(D)$,
choosing a point in its preimage $q^{-1}(z)$ only consists of fixing a flag at each point of $x_i$ with $i>0$, preserved by the Higgs field --- since the characteristic polynomial already has zero roots at all points of $D$, the Higgs field is automatically strongly parabolic.


By base change, this implies that $q_0$ is surjective as well, so we get the desired dimension upper bound
\[
\dim {\widecheck{M}^{\mathrm{par}}}_{(x,0)} \leq \dim \widecheck{M}^{\mathrm{spar}}(D)_0 = \frac{n(n-1)}{2}\deg (\CL) + (n^2-1)(g-1),
\]
which completes the proof. \qed


\begin{rmk}
\source{pw-conjecture-gl-n}
\notready
Combining Sections 3 and 4 proves that Theorem \ref{conj2.7} holds for a line bundle $\CL$ of sufficiently large degree. Then the argument of Section \ref{Sec2.3} implies that it actually holds for any effective $\CL$.
\end{rmk}








\begin{comment}
\begin{prop}
\source{pw-conjecture-gl-n}
\notready
There exists a smooth family
\[
\pi_\CM : \CM \to T
\]
together with a sheaf $\CF_\CM$ on $\CS \times_T \CM$ satisfying the following peoperties.
\begin{enumerate}
    \item[(a)] The restriction of $(\CS \times_T \CM, \CF_\CM)$ to the central fiber
    \[
    \left( \pi_\CS \times_T \pi_\CM\right)^{-1}(0) \subset \CS \times_T \CM
    \]
    is $\left(T^*C \times \CM_{\mathrm{Dol}}, \CF_n \right)$ where $\CF_n$ is a universal family of 1-dimensional sheaves.
    \item[(b)] The restriction of $(\CS \times_T \CM, \CF_\CM)$ to
    \[
    \left( \pi_\CS \times_T \pi_\CM\right)^{-1}(T^\circ) \subset \CS \times_T \CM
    \]
    is $\left((A\times \CM_{\beta,A}) \times T^\circ, \CF_\beta \boxtimes \CO_{T^\circ}\right)$.
\end{enumerate}
\end{prop}


We consider
\[
\CS \subset \bar{\CS} = \mathrm{Bl}_{C\times 0}(A \times \BP^1).
\]
The 3-fold $\bar{\CS}$ is nonsingular and projective. Let $\bar{\beta} \in H_2(\bar{\CS}, \BZ)$ be the strict transform of the curve class $n[C \times \mathrm{pt}]$ on $A \times \BP^1$. Assume that $\bar{\CM}$ is the moduli space of 1-dimensional Gieseker-stable sheaves $\CF$ (with respect to a polarization) on $\bar{\CS}$ satisfying that
\[
 [\mathrm{supp}(\CF)] = \bar{\beta}, \quad \chi(\CF) = 1.
\]
By \cite[Theorem 4.6.5]{HL}, there exists a universal family on $\bar{\CS} \times \bar{\CM}$, which is denoted by $\CF_{\bar{\CM}}$. Since $\bar{\beta}$ is a curve class supported on fibers of
\begin{equation}\label{123}
\source{pw-conjecture-gl-n}
\bar{\CS} = \mathrm{Bl}_{C\times 0}(A \times \BP^1)\rightarrow \BP^1,
\end{equation}
every $\CF \in \bar{\CM}$ is supported on a closed fiber of (\ref{123}) by the stability of $\CF$. Hence the universal family $\CF_{\bar{\CM}}$ is supported on the sub-scheme
\[
\bar{\CS} \times _{\BP^1} \bar{\CM} \subset \bar{\CS} \times \bar{\CM}.
\]

Now we define $\CM$ to be the sub-moduli space
\begin{equation}\label{CM}
\source{pw-conjecture-gl-n}
\CM \subset \bar{\CM}
\end{equation}
parametrizing 1-dimensional sheaves in $\bar{\CM}$ which are supported on $\CS \subset \bar{\CS}$. The variety $\CM$ admits a morphism
\[
\pi_\CM: \CM \to T
\]
whose closed fiber over $t\in T$ is the moduli space of sheaves in $\CM$ which are supported in the closed fiber
\[
\pi_\CS^{-1}(t) \subset \CS.
\]
The restriction of $\CF_{\bar{\CM}}$ gives a universal family $\CF'_\CM$ on $\CS \times_T \CM$, which satisfies the condition (a). Moreover, its restriction over $T^\circ \subset T$ is
\[
\CF'_\beta \boxtimes \CO_{T^\circ} \in \Coh\left((A\times  \CM_{\beta,A})\times T^\circ \right)
\]
for \emph{some} universal family $\CF'_{\beta}$ on $A\times \CM_{\beta,A}$. In particular, there exists a line bundle $\CL$ on $\CM$ such that
\[
\CF_\beta = \CF'_\beta \otimes p_\CM ^\ast \CL
\]
where $p_\CM: A \times \CM_{\beta,A} \to \CM_{\beta,A}$ is the projection.

By spreading, we obtain a line bundle $\tilde{\CL}$ on $\CM$ whose restriction over $T^\circ$ is $\CL$. Let $\tilde{p}_{\CM}: \CS \times_T\CM \to \CM$ be the natural projection. Then the variety $\CM$ together and the sheaf
\[
\CF_{\CM} = \CF'_\CM \otimes \tilde{p}_{\CM}^* \tilde{\CL}
\]
satisfies both the conditions (a) and (b). Not complete yet...
\end{comment}















\subsection{Fourier--Mukai transforms}\label{Sec2.5} Sections 2.5 to 2.7 are devoted to proving Theorem \ref{taut_abelian}. Let $E$ and $E'$ be two very general elliptic curves. We first treat the special case
\source{pw-conjecture-gl-n}
\begin{equation}\label{special}
\source{pw-conjecture-gl-n}
A= E \times E', \quad \beta = \sigma + nf
\end{equation}
where $\sigma = [\mathrm{pt} \times E']$ and $f = [E \times \mathrm{pt}]$ are the classes of a section and a fiber with respect to the elliptic fibration $p: A \to E'$.

By \cite[Theorem 3.15]{Yo}, we have an isomorphism of the moduli spaces
\[
\CM_{\beta,A} \cong M_{v_n, A}(= A^{[n]} \times \widehat{A}),\quad v_n = (1,0,-n),
\]
induced by a relative Fourier--Mukai transform with respect to the elliptic fibration $p: A \to E'$.

The support of a sheaf $\CF \in \CM_{\beta, A}$ is the union of a section and several fibers (with multiplicities). Hence the Hilbert--Chow morphism (\ref{Hilb-Ch}) is exactly the morphism
\begin{equation}\label{HC2}
\source{pw-conjecture-gl-n}
p_n\times q: A^{[n]} \times \widehat{A} \to E'^{(n)} \times E
\end{equation}
where we use $\widehat{A}\cong A = E \times E'$ and $q$ is the projection to the first factor.  Under the identification
\[
A \times A = (E \times E) \times (E'\times E'),
\]
the kernel for the relative Fourier--Mukai transform associated with the elliptic fibration $p: A \to E'$ is given by
\begin{equation}\label{rel_P}
\source{pw-conjecture-gl-n}
\CP'_E = \CP_E \boxtimes \CO_{\Delta_{E'}}.
\end{equation}
with
\[
\CP_E = \CO_{E\times E}(\Delta_E - \mathrm{pt}\times E - E\times \mathrm{pt})
\]
the normalized Poincar\'e line bundle. Recall the universal family $\CE_n$ on $A \times M_{v_n,A}$ defined in (\ref{univ_ideal}). This provides us a universal family on $A \times \CM_{\beta,A}$ of 1-dimensional Gieseker-stable sheaves,
\begin{equation}\label{univ2}
\source{pw-conjecture-gl-n}
\CF_\beta = {R\pi_{13}}_\ast \left( \pi_{12}^\ast (\pi_1^*N\otimes{\CP'_{E}}) \otimes^{\BL} \pi_{23}^\ast \CE_n \right)
\end{equation}
where $\pi_i, \pi_{ij}$ are the projection from $A \times A \times M_{v_n,A}$ to the corresponding factors, and $N \in p^* \mathrm{Pic}(E')\subset \mathrm{Pic}(A)$.\footnote{Twisting by a line bundle $N$ pulled back from $E'$ along the morphism $p: A \to {E'}$ is to ensure that the (shifted) image of $\CO_A$ with respect to the relative Fourier--Mukai transform has Euler characteristic 1; see \cite[Section 3.2]{Yo}} We define the normalized universal family
\[
\CF^{n}_\beta = {R\pi_{13}}_\ast \left( \pi_{12}^\ast {\CP'_{E}} \otimes^{\BL} \pi_{23}^\ast \CE_n \right) = \CF_\beta \otimes \pi_A^\ast N^{-1}.
\]

The universal family (\ref{univ2}) serves as the one in Theorem \ref{taut_abelian}. We construct in the following a decomposition
\begin{equation}\label{splitting1}
\source{pw-conjecture-gl-n}
H^\ast(\CM_{\beta, A}, \BQ) = \bigoplus_{i,d} \widetilde{G}_iH^\CL(\CM_{\beta, A}, \BQ)
\end{equation}
which serves as the splitting in Theorem \ref{taut_abelian}.

A canonical splitting of the perverse filtration associated with
\[
q: \widehat{A} =A \to E
\]
is given by
\[
G'_kH^\CL(\widehat{A}, \BQ) = \bigoplus_{i+k =d} H^i(E, \BQ)\otimes H^k(E',\BQ).
\]
We define (\ref{splitting1}) as the decomposition
\begin{equation}\label{splitting2}
\source{pw-conjecture-gl-n}
\widetilde{G}_kH^\ast(\CM_{\beta, A}, \BQ) = \bigoplus_{i+j=k}  \widetilde{G}_i H^\ast(A^{[n]}, \BQ) \otimes G'_jH^\ast(\widehat{A},\BQ),
\end{equation}
where
$\widetilde{G}_\ast H^\ast(A^{[n]}, \BQ)$ was defined in Section \ref{Section1.5}.
Since $\widetilde{G}_\ast H^\ast(A^{[n]}, \BQ)$ splits the perverse filtration associated with
\[
p_n:  A^{[n]} \to E'^{(n)},
\]
the filtration (\ref{splitting2}) splits the perverse filtration associated with the Hilbert--Chow morphism (\ref{Hilb-Ch}).

The following theorem verifies Theorem \ref{taut_abelian} for the special case (\ref{special}), where we can choose the twisting class to be
\[
\alpha = - \pi_A^* c_1(N) \in H^2(A\times \CM_{\beta,A}, \BQ).
\]

\begin{thm}
\source{pw-conjecture-gl-n}
\notready\label{Step1}
The decomposition (\ref{splitting2}) is multiplicative and is compatible with the normalized tautological classes; that is, we have
\begin{equation}\label{1}
\source{pw-conjecture-gl-n}
\int_\gamma \mathrm{ch}_k(\CF^n_\beta) \in \widetilde{G}_{k-1} H^\ast(\CM_{\beta,A}, \BQ),\quad \forall~\gamma \in H^\ast(A, \BQ).
\end{equation}
\end{thm}

\begin{proof}
The multiplicativity of (\ref{splitting2}) follows directly from Theorem \ref{taut_Hilb} (b). It suffices to prove (\ref{1}).

We consider the multiplicative splitting
\begin{equation}\label{eqn27}
\source{pw-conjecture-gl-n}
\widetilde{G}_kH^\ast(A \times \CM_{\beta,A} ,\BQ) = H^\ast(A, \BQ) \otimes \widetilde{G}_kH^\ast(\CM_{\beta,A}, \BQ)
\end{equation}
on the total space $A \times \CM_{\beta,A}$. This decomposition splits the perverse filtration associated with
\[
\mathrm{id}\times \pi_\beta: A \times \CM_{\beta, A} \to A \times B.
\]
The statement (\ref{1}) is equivalent to the following:
\begin{equation}\label{eqn28}
\source{pw-conjecture-gl-n}
\mathrm{ch}(\CF^n_\beta) \in \bigoplus_{k\geq 1} \widetilde{G}_{k-1} H^{2k}(A\times \CM_{\beta, A}, \BQ).
\end{equation}

For a nonsingular projective variety $X$ with a decomposition $G_*H^*(X, \BQ)$ on its cohomology, we say that a class $\alpha \in H^*(X, \BQ)$ is \emph{balanced} with respect to this decomposition if
\[
\alpha \in \bigoplus_{k} G_{k} H^{2k}(X, \BQ).
\]
We prove (\ref{eqn28}) by the following 3 steps.

\emph{Step 1.} We consider the new decomposition
\begin{equation}\label{31}
\source{pw-conjecture-gl-n}
\begin{split}
\overline{G}_k H^*(A \times \CM_{\beta, A}, \BQ)& = \bigoplus_{i+j=k} H^i(E, \BQ) \otimes H^*(E', \BQ) \otimes \widetilde{G}_jH^*(\CM_{\beta,A}, \BQ)\\
& = \bigoplus_{i+j=k} \widetilde{G}_i H^\ast(A \times A^{[n]}, \BQ) \otimes G'_jH^\ast(\widehat{A}, \BQ).
\end{split}
\end{equation}
of the cohomology $H^\ast(A\times \CM_{\beta,A}, \BQ)$. It is a multiplicative decomposition splitting the perverse filtration associated with the morphism
\[
p \times \pi_\beta: A\times \CM_{\beta,A} \to E' \times B.
\]
We first show that the class
\[
\mathrm{ch}(\CE_n) \in H^\ast(A \times M_{A, v_n}, \BQ) = H^*(A \times \CM_{\beta,A}, \BQ)
\]
is balanced with respect to the splitting (\ref{31}).


Recall the expression (\ref{univ_ideal}) of the universal family $\CE_n$. By Theorem \ref{taut_Hilb} (a), the class
\[
\mathrm{ch}(\CI_n) \in  H^{\ast}(A\times A^{[n]}, \BQ)
\]
is balanced with respect to the splitting (\ref{splitting0}). Hence via pulling back along the projection
\[
\pi_{12}: A\times \CM_{\beta,A} = A\times A^{[n]} \times \widehat{A} \to A \times A^{[n]},
\]
we obtain that the class
\[
\begin{split}
\pi_{12}^\ast \mathrm{ch}(\CI_n) =  \mathrm{ch}(\CI_n) \otimes 1_{\widehat{A}} & \in H^\ast(A \times A^{[n]}, \BQ) \otimes H^0(\widehat{A}, \BQ)\\  & \subset H^*(A\times \CM_{\beta,A}, \BQ)
\end{split}
\]
is also balanced with respect to (\ref{31}). A direct calculation shows the class $\pi_{13}^\ast c_1(\CP)$ is balanced. Hence we conclude from the multiplicativity of (\ref{31}) that the class
\[
\mathrm{ch}(\CE_n) = \pi_{12}^\ast \mathrm{ch}(\CI_n) \otimes \pi_{13}^\ast \mathrm{ch}(\CP)=  \pi_{12}^\ast \mathrm{ch}(\CI_n) \otimes \pi_{13}^\ast \mathrm{exp}(c_1(\CP))
\]
is balanced.


 \emph{Step 2.} We further consider the multiplicative decomposition
\begin{equation}\label{32}
\source{pw-conjecture-gl-n}
\widetilde{G}_k H^*(A \times A \times \CM_{\beta, A}, \BQ) = H^*(A, \BQ) \otimes \overline{G}_k H^*(A \times \CM_{\beta, A}, \BQ)
\end{equation}
of $H^*(A \times A \times \CM_{\beta, A})$, which splits the perverse filtration associated with the morphism
\[
\mathrm{id} \times p \times \pi_\beta: A \times A \times \CM_{\beta, A} \to A \times E'\times B.
\]
We show that
\begin{equation}\label{34}
\source{pw-conjecture-gl-n}
\mathrm{ch}(\pi_{12}^\ast {\CP'_{E}} \otimes^{\BL} \pi_{23}^\ast \CE_n) \in \bigoplus_{k\geq 1}\widetilde{G}_{k-1} H^{2k}(A \times A\times \CM_{\beta, A}, \BQ)
\end{equation}

In fact, by Step 1, we obtain that the class
\[
\pi_{23}^* \mathrm{ch}(\CE_n) = 1_A \otimes \mathrm{ch}(\CE_n) \in H^0(A, \BQ) \otimes H^*(A\times \CM_{\beta,A}, \BQ)
\]
is balanced with respect to (\ref{32}). The relative Poincar\'e sheaf (\ref{rel_P}) can be expressed as
\begin{equation}\label{233}
\source{pw-conjecture-gl-n}
\pi_{13}^\ast \mathrm{ch}(\CP'_E) = \pi_{13}^\ast \mathrm{ch}(\CP_E)\cup \pi_{13}^* \mathrm{ch}(\CO_{\Delta_{E'}}).
\end{equation}
The class $\pi_{12}^\ast \mathrm{ch}(\CP_E)$ is balanced via a direct calculation, and
\[
\pi_{12}^*\mathrm{ch}(\CO_{\Delta_E'}) = \pi_{12}^* c_1(\CO_{\Delta_{E'}}) \in \widetilde{G}_0 H^2(A\times A \times \CM_{\beta,A}, \BQ).
\]
Hence by (\ref{233}) and the multiplicativity of (\ref{32}), we get
\[
\pi_{12}^* \mathrm{ch}(\CP'_E)\in \bigoplus_{k\geq 1} \widetilde{G}_{k-1}H^{2k}(A\times A \times \CM_{\beta,A}).
\]
This yields (\ref{34}) since
\[
\mathrm{ch}(\pi_{12}^\ast {\CP'_{E}} \otimes^{\BL} \pi_{23}^\ast \CE_n) = \pi_{12}^*\mathrm{ch}(\CP'_E) \cup \pi_{23}^\ast \mathrm{ch}(\CE_n).
\]

 \emph{Step 3.} We finish the proof of (\ref{eqn28}).

Appying the Grothendieck--Riemann-Roch formula to the projection
\[
\pi_{13}: A \times A \times \CM_{\beta, A} \to A \times \CM_{\beta,A},
\]
we obtain that
\[
\mathrm{ch}(\CF^n_\beta) = {\pi_{13}}_\ast \mathrm{ch}(\pi_{12}^\ast {\CP'_{E}} \otimes^{\BL} \pi_{23}^\ast \CE_n).
\]
Equivalently, the class $\mathrm{ch}(\CF^n_\beta)$ corresponds to the K\"unneth factor of the class
\[
\mathrm{ch}(\pi_{12}^\ast {\CP'_{E}} \otimes^{\BL} \pi_{23}^\ast \CE_n)
\]
in the subspace
\[
H^4(A, \BQ) \otimes H^*(A \times \CM_{\beta,A}, \BQ) \subset H^*(A\times A \times \CM_{\beta, A}, \BQ).
\]
Here $H^4(A, \BQ)$ is spanned by the point class in the second factor of the product $A \times A \times \CM_{\beta, A}$. Hence (\ref{eqn28}) follows from (\ref{34}).
\end{proof}


\subsection{Perverse filtrations and $H^2(\CM_{\beta,A}, \BQ)$}
In this section, we assume that $A$ is any abelian surface with $\beta$ an ample curve class satisfying \[\beta^2= 2n \geq 4.\] Our goal is to prove a variant of Proposition \ref{prop1.1} for $\CM_{\beta,A}$.  More precisely, we introduce a \emph{canonical} class
 \begin{equation}\label{ample_base}
\source{pw-conjecture-gl-n}
 L_\beta \in H^2(\CM_{\beta,A} ,\BQ)
 \end{equation}
to characterize the perverse filtration $P_{\ast} H^*(\CM_{\beta, A})$ associated with the morphism
\[
\pi_\beta: \CM_{\beta, A} \to B.
\]


Let $\CF_0 \in \CM_{\beta, A}$ be any given point. We consider the morphisms
\[
\mathrm{det}: \CM_{\beta,A} \to \widehat{A}, \quad \CF \mapsto \mathrm{det}(\CF)\otimes \mathrm{det}(\CF_0)^\vee
\]
and
\[
\widehat{\mathrm{det}}: \CM_{\beta,A} \to A, \quad \CF \mapsto \mathrm{det}(\phi_\CP(\CF)) \otimes \mathrm{det}(\phi_\CP(\CF))^\vee.
\]
Here $\phi_\CP : D^b\Coh(A) \to D^b\Coh(\widehat{A})$ is the Fourier--Mukai transform induced by the normalized Poincar\'e line bundle $\CP$ on $A \times \widehat{A}$. By \cite{Yo}, the Albanese map for $\CM_{\beta,A}$ with respect to the reference point $\CF_0$ is
\[
\mathrm{alb} = \widehat{\mathrm{det}} \times \mathrm{det}: \CM_{\beta,A} \to A \times \widehat{A}.
\]
The closed fiber over the origin $0 \in A\times \widehat{A}$, denoted by
\[
K_{\beta,A} = \mathrm{alb}^{-1}(0) \subset \CM_{\beta, A},
\]
is a holomorphic symplectic vairiety of Kummer type. It parametrizes 1-dimensional sheaves $\CF \in \CM_{\beta,A}$ satisfying
\[
 \mathrm{det}(\CF)= \mathrm{det}(\CF_0) \in \widehat{A},\quad   \mathrm{det}(\phi_\CP(\CF)) = \mathrm{det}(\phi_\CP(\CF)) \in \widehat{\widehat{A}} = A.
\]
The restriction of (\ref{Hilb-Ch}) to the subvariety $K_{\beta, A}$ is a Lagrangian fibration
\[
\pi'_\beta: K_{\beta, A} \to  |\CO_A(\beta)| =  \BP^{n-1}.
\]
Here $ |\CO_A(\beta)|$ denotes the linear system associated with the divisor \[
\mathrm{supp}(\CF_0) \subset A.
\]

We now consider the following commutative diagram,
\begin{equation}\label{diagram1}
\source{pw-conjecture-gl-n}
    \text{[Diagram omitted; see source PDF.]}
\end{equation}
Here the horizontal arrows are defined by
\[
h(\CF, a, L) = {t_a}_* \CF \otimes  L,\quad h'(C, a)= {t_a}_\ast C
\]
with $t_a: A \xrightarrow{\simeq} A$ the translation by $a$. Since both $h$ and $h'$ are finite and surjective, the pullback morphism
\begin{equation*}
    h^*: H^k(\CM_{\beta, A}, \BQ) \to H^k(K_{\beta, A} \times A \times \widehat{A}, \BQ)
\end{equation*}
is injective. Moreover, it is an isomorphism when $k=2$ by G\"ottsche's formula \cite{Go1} for the Betti numbers of the generalized Kummer varieties.

When $n =2$, we have $\CM_{\beta,A} \cong A \times \widehat{A}$ by \cite{Yo}, and therefore $K_{\beta,A}$ is a point.

Now we consider the case $n \geq 3$. Recall the Mukai vector $v_\beta =(0,\beta ,1) \in S_A^+$. By \cite[Theorem 0.2 and (1.6)]{Yo}, the correspondence induced by the class
\[
c_1(\CF_\beta) \in H^2(A \times \CM_{\beta,A})
\]
composed with the restriction map
\[
H^2(\CM_{\beta,A}, \BQ) \to H^2(K_{\beta,A}, \BQ)
\]
gives an isometry between $v_\beta^\perp \subset S_A^+\otimes \BQ$ and $H^2(K_{\beta,A}, \BQ)$,
\[
\theta: v_\beta^\perp  \xrightarrow{\simeq}  H^2(K_{\beta,A}, \BQ).
\]
Hence the pullback $h^*$ induces an isomorphism
\begin{equation}\label{eqn39}
\source{pw-conjecture-gl-n}
\varphi: H^2(\CM_{\beta, A}, \BQ) \xrightarrow{\simeq} v_\beta^\perp \oplus H^2(A\otimes \widehat{A}, \BQ),
\end{equation}
which does not depend on the choice of $\CF_0$ by the discussion of the second factor
\[
q_{v_\beta}: H^2(\CM_{{\beta},A}, \BQ) \to H^2(A \times \widehat{A}, \BQ)
\]
in \cite[Section 8, Page 40]{Markman3}.


We consider the following classes
\[
\begin{split}
    l_\beta & = \theta\left((0,0,1)\in v_\beta^\perp\right) \in H^2(K_{\beta,A} ,\BQ), \\
    L_\beta & = \varphi^{-1} \left( l_\beta \oplus (\beta\otimes 1_{\widehat{A}}) \right) \in H^2(\CM_{\beta, A}, \BQ).
\end{split}
\]
The following proposition is the main result of this section.

\begin{prop}
\source{pw-conjecture-gl-n}
\notready\label{prop2.4}
We have
\[
P_kH^m(\CM_{\beta,A}, \BQ) = \left( \sum_{i\geq 1} \mathrm{Ker}(L_\beta^{n+k+i-m}) \cap \mathrm{Im}(L_\beta^{i-1}) \right) \cap H^m(\CM_{\beta, A}, \BQ).
\]
\end{prop}

Something to be added..









\begin{proof}
We denote by
\[
\pi_K = \pi'_\beta \times \mathrm{pr}_A : K_{\beta, A} \times A \times \widehat{A} \to \BP^{n-1}\times A
\]
the left vertical map of the diagram (\ref{diagram1}).

By \cite[Example 3.1]{Markman4}, we have
\[
l_\beta = {\pi'_\beta}^\ast \CO_{\BP^5}(1) \in H^2(K_{\beta,A}, \BQ).
\]
Hence
\[
h^*L_\beta = l_\beta \oplus (\beta\otimes 1_{\widehat{A}}) \in H^2(K_{\beta, A} \times A \times \widehat{A},\BQ)
\]
is the pullback of an ample class on the base $\BP^{n-1}\times A$ via the morphism $\pi_K$. Since $h$ and $h'$ are finite, the $t$-exactness of finite morphisms implies that a class $\alpha \in H^*(\CM_{\beta, A}, \BQ)$ lies in $P_kH^*(\CM_{\beta}, A)$ if and only if
\[
h^* \alpha \in P_kH^*(K_{\beta, A} \times A \times \widehat{A},\BQ),
\]
where the perverse filtration for $K_{\beta, A} \times A \times \widehat{A}$ is associated with the morphism $\pi_K$. Hence we obtain Proposition \ref{prop2.4} by applying Proposition \ref{prop1.1} to the morphism $\pi_K$ and the class $L = h^*L_\beta$.
\end{proof}

\subsection{Proof of Theorem \ref{taut_abelian}}
In this section, we complete the proof of Theorem \ref{taut_abelian} by combining the tools developed in the previous sections.

Let $A$ be an abelian surface and $\beta$ be an ample curve class. A \emph{perverse package} associated with the pair $(A, \beta)$ is defined to be a triple
\[
 (\CF_\beta, \alpha_\beta, L_\beta),
\]
where $\CF_\beta$ is a universal family on $A \times \CM_{\beta,A}$, $L_\beta \in H^2(\CM_{\beta,A}, \BC)$ is the class introduced in Section 2.6, and $\alpha \in H^2(A\times \CM_{\beta,A}, \BC)$ is of the type (\ref{alpha}).

An isomorphism between the perverse packages associated with $(A, \beta)$ and $(A', \beta')$ is the following:
\begin{enumerate}
    \item[(1)] A grading-preserving isomorphism of $\BC$-vector spaces
    \[
    g: H^*(A, \BC) \xrightarrow{\simeq} H^*(A', \BC).
    \]
    \item[(2)] A graded isomorphism of $\BQ$-algebras
    \[
    \phi: H^*(\CM_{\beta,A}, \BC) \xrightarrow{\simeq} H^*(\CM_{\beta',A'}, \BC).
    \]
    \item[(3)] The isomorphisms $g$ and $\phi$ satisfy that
    \[
    \begin{split}
        (g\otimes \phi) \mathrm{ch}^{\alpha_\beta} (\CF_\beta) & = \mathrm{ch}^{\alpha_{\beta'}}(\CF_{\beta'}),\\
        \phi(L_\beta) & = L_{\beta'}.
    \end{split}
    \]
\end{enumerate}
Isomorphisms of perverse packages form an equivalence relation. We call two pairs $(A, \beta)$ and $(A', \beta')$ \emph{perversely isomorphic} if there is a perverse package associated with $(A,\beta)$ isomorphic to a perverse package associated with $(A', \beta')$.

\begin{prop}
\source{pw-conjecture-gl-n}
\notready\label{prop2.5}
Assume that $(A, \beta)$ is perversely isomorphic to $(A', \beta')$. If Theorem \ref{taut_abelian} holds for $(A, \beta)$, then it also holds for $(A', \beta')$.
\end{prop}

\begin{proof}
By Proposition \ref{prop2.4}, the condition
\[
\phi(L_\beta) = L_{\beta'}
\]
implies that the perverse filtrations on $H^*(\CM_{\beta,A}, \BC)$ and $H^*(\CM_{\beta', A'}, \BC)$ are identified via $\phi$. Let
\[
H^*(\CM_{\beta,A}, \BC) = \bigoplus_{k,d} \widetilde{G}_k H^\CL(\CM_{\beta,A}, \BC)
\]
be the splitting of the perverse filtration for $\CM_{\beta, A}$ satisfying
\[
\int_\gamma \mathrm{ch}^{\alpha_\beta}_k(\CF_\beta) \in \widetilde{G}_{k-1}H^*(\CM_{\beta,A}, \BC).
\]
Then the decomposition
\[
\
\widetilde{G}_k H^\CL(\CM_{\beta',A'}, \BC) = \phi(\widetilde{G}_k H^\CL(\CM_{\beta,A}, \BC))
\]
also splits the perverse filtration for $\CM_{\beta',A'}$, and the corresponding tautological classes lie in its correct pieces by the condition (3).
\end{proof}

If $(A, \beta)$ is perversely isomorphic to $(A', \beta')$, then condition (2) implies that
\[
\beta^2 = \mathrm{dim}(\CM_{\beta,A}) -2 = \mathrm{dim}(\CM_{\beta',A'}) -2 = \beta'^2.
\]
The following theorem deduced from Theorem \ref{monodromy_main} verifies its converse.

\begin{thm}
\source{pw-conjecture-gl-n}
\notready\label{thm2.6}
Any pairs $(A, \beta)$ and $(A', \beta')$ satisfying
\[
\beta^2 = \beta'^2
\]
are perversely isomorphic.
\end{thm}

\begin{proof}

Our goal is to find a grading-preserving isomorphism of $\BC$-vector spaces
\begin{equation}\label{39}
\source{pw-conjecture-gl-n}
g: H^*(A, \BC)  \to H^*(A', \BC)
\end{equation}
satisfying the following two properties:
\begin{enumerate}
    \item[(a)] $g$ preserve the intersection pairings,
    \item[(b)] the restriction of $g$ on $S_A^+\otimes \BQ$ satisfies
    \begin{equation}\label{40}
\source{pw-conjecture-gl-n}
    \begin{split}
        g(0,0,1) & = (0,0,1),\\
        g(1,0,0) & = (1,0,0),\\
        g(0,\beta,0) & = (0, \beta',0).
    \end{split}
    \end{equation}
\end{enumerate}
and a graded ring isomorphism
 \[
 \phi : H^*(\CM_{\beta,A}, \BC) \to H^*(\CM_{\beta', A'}, \BC)
 \]
satisfying the condition (3).

We construct $(g, \phi)$ in two steps.

First, because the Mukai vectors $v_\beta$ and $v_{\beta'}$ are primitive of the same length,
we can apply \cite[Theorem 8.3]{Markman3} (which follows from Yoshioka \cite{Yo}) to find
\[ g_1: H^*(A, \BZ)  \to H^*(A', \BZ) ,\quad \phi_1: H^*(\CM_{\beta,A}, \BQ) \to H^*(\CM_{\beta', A'}, \BQ)\]
such that $g_1(v_\beta) = v_{\beta'}$ and
\begin{equation}\label{eq41}
\source{pw-conjecture-gl-n}
(g_1 \otimes \phi_1) \mathrm{ch}(\CF_\beta) = \mathrm{ch}(\CF_{\beta'})
\end{equation}
with $\CF_{\beta'}$ a universal family on $A \times \CM_{\beta,A}$.
If $\beta$ and $\beta'$ have different divisibility, then $g_1$ cannot satisfy condition b) above, so it requires further modification.

We now choose $g_2 \in \mathrm{SO}(S_{A'}^+\otimes \BC)_{v_{\beta'}}$ such that the restriction
$$g_2\circ g_1|_{S_{A}^{+}}: S_{A}^{+} \otimes \BC \rightarrow S_{A'}^{+}\otimes\BC$$
is grading-preserving and satisfies condition b).
Because we are working with complex coefficients, we can lift $g_2$ to an element of
$\mathrm{Spin}(S_{A'}^{+}\otimes \BC)_{v_{\beta'}}$ and take the associated isomorphisms
\[
\begin{split}
 g_2: &  H^*(A',\BC)  \rightarrow H^*(A',\BC)
\\
\phi_2=  \gamma_{g_2,v_{\beta'}}: &   H^*(\CM_{\beta',A'},\BC)  \rightarrow H^*(\CM_{\beta',A'},\BC).
\end{split}
\]

We deduce from Theorem \ref{monodromy_main} that
\begin{equation}\label{eq42}
\source{pw-conjecture-gl-n}
(g_2 \otimes \phi_2) \mathrm{ch}(\CF_{\beta'}) =  \mathrm{ch}^{\pi_\CM^*a}(\CF_{\beta'}), \quad a \in H^2(\CM_{\beta',A'}, \BC).
\end{equation}

Finally, we set $g = g_2 \circ g_1$ and $\phi = \phi_2\circ \phi_1$.
Then  (\ref{eq41}) and (\ref{eq42}) imply that
\[
(g\otimes \phi)\mathrm{ch}(\CF_{\beta}) = \mathrm{ch}^{\pi_\CM^*a}(\CF_{\beta'}).
\]
In particular, for any class $\alpha_\beta \in H^2(A\times \CM_{\beta,A}, \BC)$ of the type (\ref{alpha}), we have
\[
(g\otimes \phi)\mathrm{ch}^{\alpha_\beta}(\CF_{\beta}) = \mathrm{ch}^{\alpha_{\beta'}}(\CF_{\beta'}), \quad \alpha_{\beta'} = \pi_\CM^*a+ (g\otimes \phi)\alpha_\beta.
\]

Finally, under the identification (\ref{eqn39}), it follows from (\ref{40}) that the isomorphism $\phi$ induced by $g$ sends $(0,0,1) \in v_\beta^\perp$ to $(0,0,1) \in v_{\beta'}^\perp$ and sends $\beta\otimes 1_{\widehat{A}} \in H^2(A\times \widehat{A}, \BC)$ to $\beta' \otimes 1_{\widehat{A'}} \in H^2(A'\times \widehat{A'}, \BC)$. In particular, we have
\[
\phi(L_\beta) = L_{\beta'}.
\]
This completes the proof of Theorem \ref{thm2.6}.
\end{proof}

Theorem \ref{taut_abelian} follows from Theorem \ref{Step1}, Proposition \ref{prop2.5}, and Theorem \ref{thm2.6}. Later in Section \ref{Section3.5}, we will prove a strengthened version of Theorem \ref{taut_abelian} which further requires the decomposition $\tilde{G}_*H^*(\CM_{\beta,A}, \BC)$ to be induced by a construction of Deligne.

\begin{rmk}
\source{pw-conjecture-gl-n}
\notready\label{remark2.8}
It was shown in \cite{dCM} that perverse filtrations behave well under deformations. However, for the pairs $(A, \beta)$ and $(A',\beta')$ such that $\beta$ and $\beta'$ have different divisibilities in $H^2(A, \BZ)$ and $H^2(A',\BZ)$, the fibrations $\pi_\beta$ and $\pi_{\beta'}$ are not deformation equivalent.

Hence, in order to reduce Theorem \ref{taut_abelian} for any pair $(A,\beta)$ to that for a special pair
\[
A= E \times E', \quad \beta = \sigma +n f,
\]
it is essential to consider an equivalence relation weaker than deformation equivalences of fibrations $\pi_\beta$. This is our motivation to work with equivalences of perverse packages as introduced above.
\end{rmk}
